\chapter{Практика}


\section{Настройка интерфейсов роутера}

\begin{lstlisting}[language=C, caption={Пример настройки интерфейсов роутера}]

# переход в режим конфигурации
<Huawei>system-view
# выбор интерфейса
[Datacom-Router]interface GigabitEthernet 0/0/1
# назначение IP и маски
[Datacom-Router]ip address 10.0.13.1 24
# сохраняем настройки
save
\end{lstlisting}

Остальные интрерфейсы аналогично

\section{Проверка сетевой связности}

\begin{figure}[H] 
    \centering
    \includegraphics[width=0.75\textwidth]{check1.png}
    \caption{Проверка связности}
\end{figure}

Пинг проходит, значит, интерфейсы работают верно

\section{Создание LoopBack интерфейса}

\begin{lstlisting}[language=C, caption={Пример создания LoopBack интерфейса}]

interface LoopBack0
ip address 10.0.1.1 32
\end{lstlisting}

Аналогично для R2 и R3

\section{Настройка статических маршрутов}

\begin{lstlisting}[language=C, caption={Пример настройки статического маршрута}]

ip route-static 10.0.1.2 32 10.0.13.3
ip route-static 10.0.1.3 32 10.0.12.2

\end{lstlisting}

Аналогично для R2 и R3. \\

Проверка сетевой связности с LoopBack интерфейсами


\begin{figure}[H]  
    \centering
    \includegraphics[width=0.75\textwidth]{checl2.png}
    \caption{Проверка связности}
\end{figure}

\endinput